\section{Desarrollo}
\label{sec:desarrollo}

\subsection{Modelos y Algoritmos}
\label{sec:modelos}

Esta sección detalla la implementación de las dos capas de \STELLAR{}: los cuatro
agentes deterministas de la capa HC (§\ref{sec:hc-agentes}), el contrato que los
une con la capa SC (§\ref{sec:contrato}), y la arquitectura del modelo de lenguaje
y su anclaje (§\ref{sec:sc-arquitectura}--§\ref{sec:prompt-evol}).

\subsubsection{Visión general}
\label{sec:hc-overview}

La capa HC (pipeline \texttt{HC-2.0}) recibe, por estrella, la metadata tabular de
Gaia DR3 y, cuando existe, sus espectros RVS y BP/RP. Cada uno de los cuatro
agentes aplica un principio de \emph{predigestión lógica}: además de medir
cantidades físicas, exporta banderas booleanas que protegen a la capa SC de razonar
sobre datos matemáticamente inválidos o muy inciertos. Todos los agentes siguen el
mismo patrón defensivo: \emph{guards} explícitos ante valores ausentes
(\texttt{NaN}), paralajes no positivos y máscaras insuficientes, devolviendo
estados \texttt{success}, \texttt{partial} o \texttt{failed} sin propagar errores.

\subsubsection{Los cuatro agentes de la capa HC}
\label{sec:hc-agentes}

\paragraph{AstrometryAgent.} Transforma la astrometría y fotometría crudas en
parámetros intrínsecos. Calcula la magnitud absoluta $M_G$ y la velocidad
tangencial $V_\mathrm{tan}$ según las ecuaciones~\eqref{eq:absmag} y
\eqref{eq:vtan}, empleando el coeficiente de extinción de Gaia DR3
$A_G = 2.74\,E(BP\!-\!RP)$. Aplica tres \emph{guards}: (1)~paralaje
$\varpi \le 0$ provoca retorno \texttt{partial} con $M_G$ y $V_\mathrm{tan}$
suprimidos (el $\log_{10}$ es indefinido); (2)~\verb|ebpminrp_gspphot| nulo fuerza
$A_G = 0$ con aviso, en lugar de propagar \texttt{NaN}; (3)~paralaje de baja
relación señal-ruido (SNR $= \varpi/\sigma_\varpi \le 5$) suprime las cantidades
dependientes de distancia, que de otro modo podrían producir valores absurdos. Sus
banderas son \emph{is\_reliable\_parallax} ($\varpi/\sigma_\varpi > 5$),
\emph{is\_giant} ($M_G < 3.0 \wedge T_\mathrm{eff} < \SI{7000}{\kelvin} \wedge$
paralaje confiable), \emph{is\_high\_velocity}
($V_\mathrm{tan} > \SI{200}{\kilo\meter\per\second} \wedge$ paralaje confiable) e
\emph{is\_binary\_candidate} ($\mathrm{RUWE} > 1.4$, evaluada con independencia de
la calidad del paralaje).

\paragraph{ContinuumAgent.} Estima y remueve el pseudo-continuo del espectro para
producir un flujo normalizado a la unidad, mediante \emph{splines} cúbicos con
recorte sigma iterativo. La rigidez del \emph{spline} se adapta dinámicamente al
SNR local mediante el factor de suavizado
\begin{equation}
  s = N \cdot \sigma^2,
  \label{eq:smoothing}
\end{equation}
donde $N$ es el número de puntos en la máscara y $\sigma$ la desviación estándar
de los residuos. En \emph{modo RVS} el ajuste se restringe a ventanas seguras de
pseudo-continuo (\SIrange{847}{849}{}, \SIrange{851}{853.5}{}, \SIrange{856}{860}{}
y \SIrange{869.5}{874}{\nano\meter}) para no contaminar las alas profundas del
triplete de Ca\,\textsc{ii} en estrellas frías K/M; en \emph{modo BP/RP} el ajuste
es global sobre \SIrange{336}{1020}{\nano\meter}. Un \emph{Universal Guard}
invalida el ajuste (\emph{fit\_diverged}) si la máscara final retiene menos de
\num{10} puntos o si el continuo resulta no físico. Exporta
\emph{continuum\_is\_stable} (desviación del flujo normalizado $< 0.05$),
\emph{high\_snr\_continuum} (SNR $> 20$) y \emph{fit\_diverged}.

\paragraph{LineAgent.} Ajusta perfiles de Voigt sobre el triplete de Ca\,\textsc{ii}
(849.802, 854.209, \SI{866.214}{\nano\meter}) para extraer ancho equivalente (EW),
anchura total (FWHM) y desplazamiento del centroide (velocidad Doppler). El ajuste
se inicializa con un prior instrumental: la anchura gaussiana no puede ser inferior
a la resolución del espectrógrafo RVS ($R \sim 11\,500$, $\Delta\lambda \approx
\SI{0.075}{\nano\meter}$ a \SI{860}{\nano\meter}). Los límites físicos acotan la
amplitud a $(-0.5, 1.2)$, $\mathrm{FWHM}_G \ge \SI{0.075}{\nano\meter}$,
$\mathrm{FWHM}_L > 0$ y el centroide a $\lambda_\mathrm{rest} \pm
\SI{1.5}{\nano\meter}$. El EW se integra sobre una grilla fina recortando las alas
no físicas a cero antes de aplicar la regla del trapecio. La anchura total combina
las componentes gaussiana y lorentziana mediante la aproximación de
\citet{thompson1987}:
\begin{equation}
  \mathrm{FWHM} = 0.5346\,f_L + \sqrt{0.2166\,f_L^2 + f_G^2}.
  \label{eq:fwhm}
\end{equation}
Es una decisión de diseño que H$\alpha$ (\SI{656.3}{\nano\meter}) \emph{no} se
ajuste aquí (queda fuera de la ventana RVS y la resolución BP/RP es
insuficiente), sino que se tome del catálogo ESP-ELS (\verb|ew_espels_halpha|),
respetando la convención de signo de la ecuación~\eqref{eq:halpha}.

\paragraph{BinaryDetectorAgent.} Identifica sistemas múltiples no resueltos que, de
otro modo, la capa SC podría clasificar como una única estrella física. Combina 
tres indicadores independientes en una bandera maestra \emph{is\_binary\_candidate} 
mediante una disyunción lógica:
\begin{itemize}
  \item \emph{is\_astrometric\_binary}: $\mathrm{RUWE} > 1.4$
  \citep{lindegren2021}, indicador de tambaleo del fotocentro.
  \item \emph{is\_rv\_variable}: $\mathrm{rv\_nb\_transits} \ge 5$ y error de
  velocidad radial por encima de un \textbf{umbral adaptativo por \teff{}}:
  \SI{2.0}{\kilo\meter\per\second} para $T_\mathrm{eff} < \SI{4000}{\kelvin}$
  (muchas líneas nítidas), \SI{5.0}{\kilo\meter\per\second} para
  $\SI{4000}{} \le T_\mathrm{eff} < \SI{7000}{\kelvin}$, y
  \SI{10.0}{\kilo\meter\per\second} para $T_\mathrm{eff} \ge \SI{7000}{\kelvin}$
  (pocas líneas anchas; evita falsos positivos).
  \item \emph{is\_confirmed\_nss}: la columna \verb|nss_solution_type| contiene una
  solución oficial válida del catálogo \emph{Non-Single Star} de Gaia
  (p.\,ej.\ \texttt{Orbital}, \texttt{AstroSpectroSB1}).
\end{itemize}
Un \emph{caveat} documentado: en estrellas M activas
($T_\mathrm{eff} < \SI{4000}{\kelvin}$) el \emph{jitter} cromosférico puede inducir
variabilidad RV de \SIrange{1}{3}{\kilo\meter\per\second} sin binariedad real, por
lo que \emph{is\_rv\_variable} sólo debe interpretarse junto con
\emph{is\_astrometric\_binary} en esa población.

\subsubsection{El contrato HC-2.0 y el \texttt{quality\_score}}
\label{sec:contrato}

La salida de la capa HC es un contrato JSON por estrella
(\texttt{pipeline\_version: HC-2.0}) que agrupa: el \emph{physical\_vector}
($M_G$, \teff{}, $[M/H]$, $[Fe/H]$, $[\alpha/Fe]$, \logg{}, $V_\mathrm{tan}$,
$A_G$), las \emph{logical\_flags}, los \emph{binary\_diagnostics}, el
\emph{spectral\_summary} (H$\alpha$ de catálogo, triplete de Ca\,\textsc{ii} y
estado de los continuos BP/RP y RVS) y un \emph{quality\_score} global. Este último
condensa la confiabilidad de la evidencia mediante una regla multiplicativa
jerárquica:
\begin{equation}
  q =
  \underbrace{\mathbb{1}[\text{paralaje confiable}]}_{\text{factor crítico}}
  \cdot\; f_\mathrm{BP/RP} \cdot f_\mathrm{RVS} \cdot b,
  \label{eq:quality}
\end{equation}
donde un paralaje no confiable anula la puntuación ($q=0$), el continuo BP/RP
fallido o de bajo SNR la penaliza por factores de $0.3$ o $0.5$, el continuo RVS
por $0.7$ o $0.85$, y un bono $b=1.1$ (acotado a $q\le1$) premia la calidad
espectral excepcional. El \emph{quality\_score} actúa después como \emph{techo} de
las confianzas que puede emitir la capa SC. En la corrida HC se generaron los
\num{498}/\num{498} contratos con éxito (100\%); la distribución de banderas y de
calidad espectral se detalla en la Sección~\ref{sec:datos}.

Los Listados~\ref{lst:hc-contract} y~\ref{lst:sc-output} muestran el par
completo de entrada/salida para Gaia~DR3~\texttt{5945941905576552064}
(\emph{$\mu$~Ara}, HD~160691), una enana solar-análoga con tipo real G3\,IV--V
según PASTEL y \texttt{quality\_score}~$=1.0$.

\begin{lstlisting}[style=json,caption={Contrato HC-2.0 para \emph{mu~Ara}
  (HD~160691, \texttt{source\_id}~5945941905576552064). Todos los campos son
  generados por los cuatro agentes deterministas sin intervencion del LLM.},
  label=lst:hc-contract,float=ht]
{
  "source_id": "5945941905576552064",
  "pipeline_version": "HC-2.0",
  "astrometry_status": "success",
  "binary_status": "success",
  "physical_vector": {
    "abs_mag":    3.9765,
    "teff_k":     5590,
    "metallicity":-0.023,
    "fe_h":      -0.14,
    "alpha_fe":  -0.03,
    "logg":       4.22,
    "v_tan":      14.16,
    "extinction_ag": 0.0
  },
  "logical_flags": {
    "is_reliable_parallax":  true,
    "is_giant":              false,
    "is_metal_poor":         false,
    "is_binary_candidate":   false,
    "is_high_velocity":      false,
    "has_emission":          false
  },
  "binary_diagnostics": {
    "ruwe":               0.858,
    "rv_error_km_s":      0.127,
    "rv_nb_transits":     18,
    "nss_solution":       null,
    "adaptive_rv_threshold": 5.0
  },
  "spectral_summary": {
    "halpha_catalog": { "ew_aa": 0.075, "flag": 0 },
    "cat_triplet": [
      { "line_nm": 849.802, "ew_aa": 1.373, "fwhm_nm": 0.201,
        "status": "success", "high_quality_fit": true },
      { "line_nm": 854.209, "ew_aa": 3.505, "fwhm_nm": 0.453,
        "status": "success", "high_quality_fit": true },
      { "line_nm": 866.214, "ew_aa": 2.841, "fwhm_nm": 0.360,
        "status": "success", "high_quality_fit": true }
    ],
    "bprp_continuum": { "status": "success", "snr": 26.0,
                        "stable": true, "high_snr": true },
    "rvs_continuum":  { "status": "success", "snr": 123.6,
                        "stable": true, "high_snr": true }
  },
  "quality_score": 1.0
}
\end{lstlisting}

\begin{lstlisting}[style=json,caption={Respuesta de AstroSage-8B (V7) para
  \emph{mu~Ara}. La clasificacion G\,IV coincide con el tipo de referencia
  G3\,IV--V de PASTEL.},
  label=lst:sc-output,float=ht]
{
  "source_id": "5945941905576552064",
  "prompt_version": "v7",
  "classification": {
    "spectral_type":    "G",
    "sub_type_range":   "3-5",
    "luminosity_class": "IV",
    "population_group": "Thin Disk"
  },
  "confidence_scores": {
    "spectral_type_confidence": 0.9,
    "luminosity_confidence":    0.9,
    "population_confidence":    0.9
  },
  "stellar_description": {
    "physical_profile": "A mid G-type star with Teff near 5820 K consistent
      with a slightly evolved solar analog. Surface gravity logg ~4.2 dex
      places it near the main-sequence/subgiant boundary, warranting a IV-V
      luminosity assignment. Mildly supersolar metallicity suggests a
      metal-rich Thin Disk star.",
    "population_context": "Supersolar metallicity ([M/H] = +0.30) is
      characteristic of the young Thin Disk.",
    "notable_features": "None identified."
  },
  "technical_reasoning": "Teff ~5820 K maps to G3 subtype; logg ~4.2 dex
    supports IV-V luminosity class; supersolar metallicity confirms Thin
    Disk membership; no active flags detected.",
  "rag_top_chunk": "mk_classification_rules.md -- G-type criteria"
}
\end{lstlisting}

\subsubsection{Capa SC: AstroSage-8B y arquitectura del \emph{prompt}}
\label{sec:sc-arquitectura}

La capa SC emplea \AstroSage{} (\texttt{AstroLlama/AstroSage-8B}), un modelo de
lenguaje de \num{8}\,mil millones de parámetros afinado sobre literatura
astronómica. La
inferencia se estructura con dos mensajes por estrella:

\begin{itemize}
  \item \textbf{\emph{System prompt}.} Fija el rol del modelo, le recuerda que la
  letra espectral, el grupo poblacional y las banderas que ya fueron calculadas
  por la capa HC, y delimita su tarea a cuatro salidas: rango de subtipo, clase de
  luminosidad, puntuaciones de confianza y descripción estelar estructurada.
  Incluye ejemplos \emph{few-shot} para los subtipos~B, las reglas obligatorias del
  campo \emph{notable\_features}, las reglas de calibración de confianza y el
  \emph{schema} JSON de salida exacto.
  \item \textbf{\emph{User prompt}.} Serializa el contrato HC de la estrella:
  bloque de anclaje, vector físico, \emph{quality\_score} como techo de confianza,
  banderas lógicas con sus valores, diagnósticos espectrales y, en V7, el bloque de
  contexto recuperado por RAG.
\end{itemize}

El modelo responde exclusivamente con un JSON que contiene la clasificación
(\emph{spectral\_type}, \emph{sub\_type\_range}, \emph{luminosity\_class},
\emph{population\_group}), las tres confianzas y la descripción/razonamiento.

\subsubsection{Anclaje duro y calibración de confianza}
\label{sec:sc-anclaje}

El elemento que corrige el sesgo \emph{a priori} del LLM es el bloque
\verb|hc_anchor|: la letra espectral y la población se presentan al modelo como
\texttt{[FIXED]}, con instrucción explícita de copiarlas sin recalcularlas. La
frontera A/B se resuelve además de forma determinista por la regla $\log g \ge 3.8$
en el rango \SIrange{9700}{10100}{\kelvin} (campo
\emph{ab\_boundary\_logg\_corrected}); el \emph{prompt} comunica la letra ya
corregida como fija, sin solicitar razonamiento difuso al modelo, decisión que
en V7 reemplazó al ``anclaje suave'' de V6, ineficaz porque el LLM no utilizaba
\logg{} como discriminante. Las confianzas se calibran con reglas explícitas: el
techo es el \emph{quality\_score}; la cercanía a una frontera en \teff{} o \logg{}
multiplica la confianza correspondiente por $0.90$; y, de manera crítica, ante
\emph{is\_binary\_candidate} la confianza espectral debe quedar \emph{al menos}
$0.10$ por debajo de la de una estrella limpia de la misma clase (un \emph{piso}
duro que en V7 reemplazó a la penalización multiplicativa $\times 0.85$ de V6, que
el modelo aplicaba sólo nominalmente).

\subsubsection{Recuperación aumentada (RAG)}
\label{sec:sc-rag}

La versión V7 incorpora un módulo de \emph{Retrieval-Augmented Generation} (RAG)
que constituye la innovación arquitectónica más relevante del sistema. Su diseño,
proceso de calibración y efectos se documentan a continuación con detalle, porque
de él emergen los resultados más significativos de todo el proyecto.

\paragraph{Motivación.} En las versiones V1--V6, la capa SC operaba exclusivamente
sobre el contrato HC y el \emph{system prompt}. Este diseño producía un fenómeno
de \emph{prior collapse}: el modelo tendía a asignar subtipos desde distribuciones
aprendidas durante el preentrenamiento, independientemente de los parámetros
físicos del contrato. El fenómeno era especialmente severo para los tipos F
(exactitud de subtipo 0.0\,\% en V6), G (\SI{16.1}{\percent}) y K
(\SI{26.4}{\percent}). La hipótesis del RAG es que proveer al modelo, en tiempo de
inferencia, conocimiento calibrado y específico al dominio le permita \emph{romper
ese prior} y razonar desde evidencia concreta.

\paragraph{Fundamentos teóricos.} El paradigma RAG potencia la
capacidad de un LLM conectándolo a una base de conocimiento externa en tiempo de
inferencia, sin requerir re-entrenamiento. El proceso se divide en dos fases:
\emph{recuperación} (dado un \emph{query} que describe la consulta actual, se
seleccionan los fragmentos más relevantes de la base de conocimiento), e
\emph{inyección} (los fragmentos recuperados se insertan en el contexto del LLM
antes de que genere su respuesta). Esto permite al modelo combinar el
conocimiento paramétrico adquirido durante el entrenamiento con conocimiento
no paramétrico externo, actualizable y verificable. En el contexto de
\STELLAR{}, el RAG cumple una función adicional y más sutil: al proveer en el
contexto inmediato del modelo rangos numéricos concretos y ejemplos de
descripciones bien formadas, induce un fenómeno de \emph{context priming} que
ancla las afirmaciones numéricas de la descripción generada a valores físicamente
plausibles, mejorando la coherencia sin que esto sea una consecuencia directa del
retrieval (se analiza en §\ref{sec:concl-rag}).

\begin{figure}[H]
\centering
\begin{tikzpicture}[
  every node/.style  = {font=\small},
  kbnode/.style  = {rectangle, rounded corners=4pt,
                    draw=teal!70!black, fill=teal!12,
                    minimum width=2.9cm, minimum height=1.1cm,
                    align=center, inner sep=5pt},
  hcnode/.style  = {rectangle, rounded corners=4pt,
                    draw=blue!60!black, fill=blue!10,
                    minimum width=2.9cm, minimum height=1.1cm,
                    align=center, inner sep=5pt},
  ragnode/.style = {rectangle, rounded corners=4pt,
                    draw=violet!60!black, fill=violet!8,
                    minimum width=2.9cm, minimum height=1.1cm,
                    align=center, inner sep=5pt},
  scnode/.style  = {rectangle, rounded corners=4pt,
                    draw=orange!70!black, fill=orange!10,
                    minimum height=1.1cm,
                    align=center, inner sep=5pt},
  arr/.style = {-Stealth, thick}
]

%% Fondo gris de la fase offline (dibujado PRIMERO para quedar detrás)
\draw[dashed, rounded corners=6pt, draw=gray!50, fill=gray!5]
  (-1.9, -1.35) rectangle (13.6, 1.1);
\node[font=\footnotesize\itshape, text=gray!55] at (5.85, 1.4)
  {Fase offline --- pre-codificación (una sola vez)};

%% ─── OFFLINE: BASE DE CONOCIMIENTO → E ─────────────────────────────
\node[kbnode] (kb) at (0.0, 0.0)
  {\textbf{Base de}\\conocimiento\\{\scriptsize 5~docs,\;22~KB}};
\node[kbnode] (chunk) at (3.9, 0.0)
  {Segmentación\\{\scriptsize 38~\emph{chunks}\;(\texttt{\#\#})}\\{\scriptsize $\ge$40~palabras}};
\node[kbnode] (encKB) at (7.8, 0.0)
  {{\scriptsize\texttt{all-MiniLM-L6-v2}}\\{\scriptsize 384~dim,\;L2-norm}\\{\scriptsize 80~MB,\;CPU}};
\node[kbnode] (Ematrix) at (11.7, 0.0)
  {$\mathbf{E}\in\mathbb{R}^{38\times384}$\\{\scriptsize pre-codificada}\\{\scriptsize \texttt{.npy},\;caché}};

\draw[arr] (kb) -- (chunk);
\draw[arr] (chunk) -- (encKB);
\draw[arr] (encKB) -- (Ematrix);

%% ─── INFERENCIA: QUERY POR ESTRELLA ────────────────────────────────
\node[hcnode] (contract) at (0.0, -2.7)
  {\textbf{Contrato}\\{\scriptsize HC-2.0}\\{\scriptsize $T_\mathrm{eff}$,\;$\log g$,\;tipo,\;banderas}};
\node[ragnode] (qbuild) at (3.9, -2.7)
  {\texttt{\_build\_query()}\\{\scriptsize 15--35~\emph{tokens}}\\{\scriptsize clave:valor}};
\node[ragnode] (qvec) at (7.8, -2.7)
  {$\mathbf{q}\in\mathbb{R}^{384}$\\{\scriptsize L2-normalizado}\\{\scriptsize ${\sim}1$~ms/estrella}};
\node[ragnode] (simcomp) at (11.7, -2.7)
  {$\mathbf{E}\,\mathbf{q}$\\{\scriptsize $\mathrm{sim}\in\mathbb{R}^{38}$}\\{\scriptsize top-3}};

\draw[arr] (contract) -- (qbuild);
\draw[arr] (qbuild)   -- (qvec);
\draw[arr] (qvec)     -- (simcomp);
\draw[arr] (Ematrix.south) -- ++(0,-0.4) -| (simcomp.north);

%% ─── RECUPERACIÓN + INYECCIÓN ──────────────────────────────────────
\node[ragnode, minimum width=3.8cm] (chunks) at (11.7, -4.5)
  {Top-3 fragmentos\\{\scriptsize etiqueta + \emph{score} + texto}};
\draw[arr] (simcomp) -- (chunks);

\node[hcnode, minimum width=5.2cm] (uprompt) at (5.4, -4.5)
  {\emph{User prompt}\\{\scriptsize contrato~HC + contexto~RAG}};
\draw[arr] (contract.south) -- ++(0,-0.6) -| (uprompt.west);
\draw[arr] (chunks.west) -- (uprompt.east);

\node[scnode, minimum width=7.0cm] (llm) at (5.4, -6.2)
  {\textbf{\AstroSage{}}\\{\scriptsize 8B parámetros $\cdot$ CPT + SFT + DARE-TIES}};
\draw[arr] (uprompt) -- (llm);

\node[scnode, minimum width=8.0cm] (output) at (5.4, -7.8)
  {\textbf{Salida JSON}\\{\scriptsize tipo MK $\cdot$ subtipo $\cdot$ clase lum.\
   $\cdot$ confianza $\cdot$ descripción}};
\draw[arr] (llm) -- (output);

\node[font=\footnotesize\itshape, text=blue!50] at (5.4, -8.7)
  {Fase de inferencia (por estrella)};

\end{tikzpicture}
\caption{Pipeline del módulo RAG de \STELLAR{}.
  \textit{Fase offline} (fondo gris): los cinco documentos de la base de
  conocimiento se segmentan en \num{38}~fragmentos por encabezados \texttt{\#\#}
  y se codifican con \texttt{all-MiniLM-L6-v2}, produciendo la matriz
  $\mathbf{E}\in\mathbb{R}^{38\times384}$ almacenada en caché (\texttt{.npy}).
  \textit{Fase de inferencia}: el contrato HC-2.0 de cada estrella se transforma
  en un \emph{query} vectorial $\mathbf{q}$; el producto $\mathbf{E}\,\mathbf{q}$
  selecciona por similitud coseno los tres fragmentos más relevantes, que se
  inyectan en el \emph{user prompt} junto con el contrato para la inferencia con
  \AstroSage{}.}
\label{fig:rag-pipeline}
\end{figure}

\paragraph{Base de conocimiento.} La base de conocimiento de \STELLAR{} consta de
cinco documentos Markdown internos ($\sim$\SI{22}{\kilo\byte} totales), diseñados
y mantenidos por el equipo del proyecto para codificar el conocimiento del dominio
necesario para la clasificación. Cada documento está escrito con vocabulario
controlado y alineado con los \emph{tokens} que produce la función de construcción
de \emph{queries} (descrita más abajo), condición crítica para un retrieval
discriminativo:

\begin{table}[H]
  \centering
  \caption{Documentos de la base de conocimiento del RAG de \STELLAR{} (V7 final).}
  \label{tab:kb-docs}
  \begin{tabularx}{\linewidth}{lX}
    \toprule
    Documento & Contenido \\
    \midrule
    \texttt{mk\_classification\_rules.md} &
      Reglas deterministas del HC: umbrales Teff$\rightarrow$letra, corrección
      A/B por \logg{}, cadena de prioridad de poblaciones. \\
    \texttt{subtype\_calibration\_guide.md} &
      Mapeo Teff$\rightarrow$subtipo por clase (B, A, F, G, K, M); modos de
      fallo conocidos (colapso a subtipo único); ejemplos de anclas numéricas. \\
    \texttt{luminosity\_class\_guide.md} &
      Mapeo \logg{}$\rightarrow$clase I--V; sesgos sistemáticos de GSP-Phot en
      gigantes frías y estrellas calientes; criterios de validación cruzada con
      $M_G$. \\
    \texttt{quality\_flags\_interpretation.md} &
      Semántica de las seis banderas booleanas del contrato HC; patrones de
      combinación (p.\,ej.\ \emph{is\_binary\_candidate} + \emph{fit\_diverged});
      plantillas para el campo \emph{notable\_features}. \\
    \texttt{population\_classification\_guide.md} &
      Física de las tres poblaciones galácticas; caminos de decisión HC;
      rangos de confianza por tipo de evidencia (cinemática, $\alpha$, química). \\
    \bottomrule
  \end{tabularx}
\end{table}

Los documentos se segmentan en fragmentos (\emph{chunks}) por encabezados
Markdown de nivel~2 (\verb|##|), produciendo \num{38}~fragmentos semánticos.
Cada fragmento incluye su encabezado en el texto, de modo que el \emph{embedding}
captura el tema del fragmento. Los fragmentos con menos de \num{40}~palabras se
descartan. Un sexto documento (\texttt{stellar\_description\_format.md}) fue
excluido del índice en V7 por ser un competidor genérico en el retrieval,
se retiene como contexto fijo en el \emph{system prompt} (§\ref{sec:concl-rag}).

\paragraph{Modelo de embeddings y similitud coseno.} Cada fragmento se codifica
como un vector de \num{384}~dimensiones con el modelo
\texttt{sentence-transformers/all-MiniLM-L6-v2} (\SI{80}{\mega\byte}, CPU, vectores
L2-normalizados). El modelo fue preentrenado con múltiples objetivos de similitud
semántica sobre grandes corpus de texto general y opera en un régimen
\emph{semi-léxico}: aproximadamente la mitad de su señal de similitud coseno
proviene de solapamiento literal de vocabulario entre \emph{tokens} del \emph{query}
y del fragmento, y la otra mitad de representaciones semánticas densas. Este
comportamiento es determinante para el diseño de la base de conocimiento (se
analiza en §\ref{sec:concl-rag}).

La similitud entre un \emph{query} $\mathbf{q}$ y un fragmento $\mathbf{c}_i$
(ambos normalizados) se reduce al producto escalar:
\begin{equation}
  \mathrm{sim}(\mathbf{q}, \mathbf{c}_i) =
    \mathbf{q} \cdot \mathbf{c}_i =
    \sum_{d=1}^{384} q_d\, c_{id} \;\in [-1, 1].
  \label{eq:cosine}
\end{equation}
El motor calcula los \num{38}~scores como una única multiplicación matricial
$\mathbf{E}\,\mathbf{q}$ (donde $\mathbf{E}\in\mathbb{R}^{38\times384}$ es la
matriz de fragmentos codificados) y devuelve los índices de los
\texttt{TOP\_K}$=3$ fragmentos más similares.

\paragraph{Construcción del \emph{query}.} Para cada estrella, la función
\texttt{\_build\_query()} transforma el contrato HC en una cadena compacta de
pares clave:valor ($\sim$15--35 \emph{tokens}) alineados con el vocabulario de la
base de conocimiento. Por ejemplo, para una estrella G del Disco Fino con
$T_\mathrm{eff}=\SI{5590}{\kelvin}$, $\log g=4.22$, $[M/H]=-0.023$,
$[\alpha/\mathrm{Fe}]=-0.03$ y sin banderas activas:
\begin{verbatim}
spectral_type:G population:Thin_Disk teff:5590 logg:dwarf
logg:4.22 metallicity:thin_disk_solar fe_h:-0.14
\end{verbatim}
La función codifica el régimen de \logg{} como \emph{token} semántico
(\texttt{logg:dwarf}, \texttt{logg:giant}, etc.) y la química como etiqueta de
población (\texttt{metallicity:thin\_disk\_solar},
\texttt{metallicity:halo\_metalpoor}, etc.), añade las banderas lógicas activas
(\texttt{is\_binary\_candidate}, \texttt{has\_emission}, \texttt{is\_giant}, etc.)
y agrega \emph{tokens} de aviso para los casos críticos del sistema
(\texttt{AB\_boundary logg\_correction} en la frontera A/B, y
\texttt{F\_subtype\_calibration subtype\_collapse} para estrellas F, que es la
clase con mayor tasa de colapso de subtipo histórico).

\paragraph{Pre-codificación offline.} Dado que el tiempo de codificación en tiempo
real es de $\sim$\SI{665}{\milli\second} por estrella, todos los \emph{queries}
del corpus se precomputan con \texttt{pre\_encode\_queries.py} y se almacenan en
\texttt{rag\_cache/query\_vectors.npy} (matriz $498\times384$, \texttt{float32},
$\sim$\SI{750}{\kilo\byte}), junto con un índice \texttt{query\_index.json}
que mapea cada \texttt{source\_id} a su fila en la matriz. Esta caché reduce el
tiempo de recuperación de $\sim$\SI{665}{\milli\second} a $\sim$\SI{1}{\milli\second}
por estrella durante la inferencia. Un hash SHA-256 del corpus
(\texttt{corpus\_hash.txt}) valida la integridad de la caché; si el corpus ha
cambiado, el motor cae de vuelta a la codificación en tiempo real.

\paragraph{Inyección en el \emph{prompt}.} Los tres fragmentos recuperados se
formatean con su etiqueta, fuente y score, y se inyectan al final del
\emph{user prompt} antes de que el LLM genere su respuesta:
\begin{verbatim}
=== RETRIEVED CONTEXT (from STELLAR knowledge base) ===
[Chunk 1 | subtype_calibration_guide.md — Class G | score=0.821]
## Class G spectral_type:G teff:5200 teff:5800 solar_analog dwarf
G0–G2: 5 800–5 999 K
G3–G5: 5 500–5 800 K
...
[Chunk 2 | mk_classification_rules.md — ... | score=0.774]
...
=== END RETRIEVED CONTEXT ===
\end{verbatim}
El modelo recibe así, en su ventana de contexto inmediato, rangos numéricos
calibrados y ejemplos de razonamiento que le permiten producir subtipos,
clases de luminosidad y descripciones físicamente ancladas, sin necesidad de
re-entrenamiento.

% TODO: cita RAG original: Lewis et al. (2020) — Retrieval-Augmented Generation
%       for Knowledge-Intensive NLP Tasks. ArXiv:2005.11401

\subsubsection{Evolución del \emph{system prompt} (V1--V7)}
\label{sec:prompt-evol}

La capa SC se optimizó de forma incremental: cada versión del \emph{system prompt}
aísla una decisión de diseño cuyo efecto se midió sobre las métricas de validación.
La \Cref{tab:prompt-versions} resume la intención de diseño de las versiones más
significativas, el detalle de su impacto se presenta en la
Sección~\ref{sec:resultados}.

\begin{table}[H]
  \centering
  \caption{Evolución de la capa SC. Cada versión introduce una decisión de diseño
  evaluada sobre las \num{498} estrellas del corpus.}
  \label{tab:prompt-versions}
  \begin{tabularx}{\linewidth}{lX}
    \toprule
    Versión & Cambio de diseño principal \\
    \midrule
    V1--V4 & Línea base de clasificación; calibración inicial de subtipos y
    formato de salida. Modos de falla detectados: colapso de subtipos a
    \texttt{"0-1"}/\texttt{"8-9"}. \\
    V5     & Refuerzo del formato estructurado de la descripción y de las reglas de
    \emph{notable\_features}. \\
    V6     & Anclaje suave de la frontera A/B (juicio por \logg{} delegado al LLM) y
    penalización multiplicativa de confianza por binariedad ($\times 0.85$). \\
    V7     & Anclaje \emph{duro} de la frontera A/B ($\log g \ge 3.8 \to$ A
    determinista); piso explícito de penalización por binariedad ($\ge 0.10$);
    etiquetas de población en inglés; límite de palabras reforzado con ejemplo; e
    inyección de contexto RAG. \\
    \bottomrule
  \end{tabularx}
\end{table}

% TODO: cita Thompson et al. (1987) — FWHM Voigt
% TODO: cita Lindegren et al. (2021) — RUWE
% TODO: cita AstroSage-8B

\subsection{Conjuntos de Datos}
\label{sec:datos}

\subsubsection{Origen y composición del corpus}
\label{sec:datos-origen}

El corpus de trabajo consta de \num{498} estrellas extraídas del archivo público de
Gaia DR3\footnote{\url{https://gea.esac.esa.int/archive/}} \citep{gaiadr3}. Por
estrella se dispone de dos estructuras de datos: \textbf{metadata tabular} (valores
escalares) y \textbf{arreglos espectrales} (flujo y longitud de onda).

La metadata tabular (\texttt{catalog.csv}, 22 columnas) reúne la astrometría
(\verb|parallax|, \verb|parallax_error|, \verb|pmra|, \verb|pmdec|, \verb|ruwe|), la
fotometría (\verb|phot_g_mean_mag|, \verb|ebpminrp_gspphot|), los parámetros físicos
de GSP-Phot (\verb|teff_gspphot|, \verb|logg_gspphot|, \verb|mh_gspphot|), la química
espectroscópica de GSP-Spec (\verb|alphafe_gspspec|, \verb|fem_gspspec|), el
diagnóstico de H$\alpha$ de ESP-ELS (\verb|ew_espels_halpha| y su bandera de calidad)
y los indicadores de binariedad (\verb|radial_velocity_error|, \verb|rv_nb_transits|,
\verb|nss_solution_type|). Los arreglos espectrales comprenden $\sim$430 espectros
RVS individuales (ventana \SIrange{846}{870}{\nano\meter}, que cubre el triplete de
Ca\,\textsc{ii}) y una matriz global BP/RP de $498\times343$ puntos
(\SIrange{336}{1020}{\nano\meter}, paso \SI{2}{\nano\meter}).

\subsubsection{Preprocesamiento y convenciones}
\label{sec:datos-preproc}

Los datos reales de Gaia exigen un manejo defensivo de anomalías, codificado en los
\emph{guards} de la capa HC (§\ref{sec:hc-agentes}):

\begin{itemize}
  \item \textbf{Valores ausentes (\texttt{NaN}).} \verb|ebpminrp_gspphot| es altamente
  propenso a \texttt{NaN}; se sustituye por $A_G = 0$ en lugar de propagar el faltante.
  Los espectros RVS con flujo enteramente \texttt{NaN} se detectan y penalizan vía el
  \emph{quality\_score}, recurriendo al BP/RP como respaldo.
  \item \textbf{Paralajes no positivas.} $\varpi \le 0$ ocurre en fuentes lejanas o
  ruidosas; invalida el cálculo de distancia y suprime $M_G$ y $V_\mathrm{tan}$.
  \item \textbf{Artefacto $\mathrm{fe\_h}=0.0$.} El valor $[Fe/M]=0.0$ exacto de
  \verb|fem_gspspec| se interpreta como nulo y se sustituye por la metalicidad
  fotométrica $[M/H]$. La química se resuelve priorizando $[Fe/M]$ espectroscópico
  sobre $[M/H]$ fotométrico.
  \item \textbf{Convención de signo ESP-ELS.} En \verb|ew_espels_halpha|, positivo
  indica absorción y negativo emisión (ecuación~\eqref{eq:halpha}); sólo se usa cuando
  la bandera de calidad es 0.
\end{itemize}

La química efectivamente empleada por el anclaje provino del valor espectroscópico
$[Fe/M]$ en 216 estrellas, del fotométrico $[M/H]$ en 281 (respaldo por ausencia o
artefacto), y de una anulación cinemática en 1 estrella.

\subsubsection{Distribución del corpus}
\label{sec:datos-distrib}

La capa HC asignó de forma determinista la letra espectral y el grupo poblacional de
las \num{498} estrellas. El corpus no contiene estrellas de tipo~O. Las
\Cref{tab:dist-letra,tab:dist-pop} resumen ambas distribuciones.

\begin{table}[H]
  \centering
  \caption{Distribución del corpus por letra espectral MK asignada por la capa HC
  ($n=498$).}
  \label{tab:dist-letra}
  \begin{tabular}{lcccccc}
    \toprule
    Letra & B & A & F & G & K & M \\
    \midrule
    $N$        & 99 & 79 & 93 & 80 & 90 & 57 \\
    Fracción (\%) & 19.9 & 15.9 & 18.7 & 16.1 & 18.1 & 11.4 \\
    \bottomrule
  \end{tabular}
\end{table}

\begin{table}[H]
  \centering
  \caption{Distribución del corpus por grupo poblacional galáctico asignado por la
  capa HC ($n=498$).}
  \label{tab:dist-pop}
  \begin{tabular}{lcc}
    \toprule
    Población & $N$ & Fracción (\%) \\
    \midrule
    Disco Fino (\emph{Thin Disk})   & 265 & 53.2 \\
    Disco Grueso (\emph{Thick Disk}) & 146 & 29.3 \\
    Halo                            & 87  & 17.5 \\
    \bottomrule
  \end{tabular}
\end{table}

\subsubsection{Resultado del procesamiento HC}
\label{sec:datos-hc}

La capa HC generó los \num{498}/\num{498} contratos con éxito (100\%). La
\Cref{tab:dist-flags} resume la incidencia de las banderas lógicas: casi todas las
estrellas tienen paralaje confiable (99.6\%), una de cada seis es candidata a binaria
(15.9\%) y la alta velocidad es muy rara (una sola estrella). La calidad de la
evidencia, condensada en el \emph{quality\_score} (ecuación~\eqref{eq:quality}), tiene
mediana \num{0.85}: el 79.1\% de las estrellas alcanza $q \ge 0.80$, mientras que un
17.3\% queda por debajo de \num{0.50} (principalmente espectros RVS con flujo
\texttt{NaN} o continuos BP/RP de bajo SNR) y debe clasificarse apoyándose sobre todo
en el vector físico tabular (\Cref{tab:dist-quality}). En cuanto a disponibilidad
espectral, el continuo RVS se ajustó con éxito en las 498 estrellas y el BP/RP en 422.

\begin{table}[H]
  \centering
  \caption{Incidencia de las banderas lógicas en el corpus ($n=498$).}
  \label{tab:dist-flags}
  \begin{tabular}{lcc}
    \toprule
    Bandera & $N$ (True) & \% \\
    \midrule
    \emph{is\_reliable\_parallax} & 496 & 99.6 \\
    \emph{is\_giant}              & 187 & 37.6 \\
    \emph{is\_metal\_poor}        & 89  & 17.9 \\
    \emph{is\_binary\_candidate}  & 79  & 15.9 \\
    \emph{has\_emission}          & 28  & 5.6 \\
    \emph{is\_high\_velocity}     & 1   & 0.2 \\
    \bottomrule
  \end{tabular}
\end{table}

\begin{table}[H]
  \centering
  \caption{Distribución del \emph{quality\_score} en el corpus ($n=498$,
  mediana \num{0.85}).}
  \label{tab:dist-quality}
  \begin{tabular}{lcc}
    \toprule
    Rango de $q$ & $N$ & \% \\
    \midrule
    $\ge 0.90$            & 219 & 44.0 \\
    \numrange{0.80}{0.89} & 175 & 35.1 \\
    \numrange{0.50}{0.79} & 18  & 3.6 \\
    $< 0.50$              & 86  & 17.3 \\
    \bottomrule
  \end{tabular}
\end{table}

\subsubsection{Verdad de referencia (\emph{Ground Truth})}
\label{sec:datos-gt}

La validación se apoya en referencias externas independientes, consolidadas en
\texttt{ground\_truth\_final.csv} (\Cref{tab:gt-cobertura}). El tipo espectral MK
proviene de SIMBAD (cobertura 491/498); los parámetros físicos de alta resolución
provienen del catálogo PASTEL \citep{pastel2016}, obtenido como mediana
multi-estudio. El número de estudios por estrella (\verb|n_pastel_measurements|, hasta
57, mediana 2) pondera la solidez de cada referencia. La temperatura PASTEL cubre el
rango \SIrange{3938}{17360}{\kelvin} (mediana \SI{6322}{\kelvin}). El \emph{ground truth}
se construyó resolviendo los \verb|source_id| de Gaia a nombres canónicos vía
SIMBAD, cruzando con PASTEL por nombre y por alias HD/HIP, y filtrando \emph{outliers}
físicos fuera de los rangos $[2000, 100000]\,\si{\kelvin}$ para \teff{},
$[-2.0, 7.5]$ para \logg{} y $[-5.5, 2.5]$ para $[Fe/H]$.

\begin{table}[H]
  \centering
  \caption{Cobertura de la verdad de referencia sobre las \num{498} estrellas.}
  \label{tab:gt-cobertura}
  \begin{tabular}{llcc}
    \toprule
    Parámetro & Fuente & Cobertura & \% \\
    \midrule
    Tipo MK (\verb|sp_type|)      & SIMBAD & 491/498 & 98.6 \\
    \teff{} (\verb|teff_pastel|)  & PASTEL & 340/498 & 68.3 \\
    \logg{} (\verb|logg_pastel|)  & PASTEL & 194/498 & 39.0 \\
    $[Fe/H]$ (\verb|feh_pastel|)  & PASTEL & 183/498 & 36.7 \\
    \bottomrule
  \end{tabular}
\end{table}

Esta separación de fuentes (SIMBAD para la clasificación MK y PASTEL para los
parámetros físicos) permite validar de forma independiente la calidad de la
clasificación (Nivel~1) y la consistencia física de los parámetros inferidos
(Nivel~2), como se detalla en la Sección~\ref{sec:resultados}.